\documentclass[12pt,a4paper]{article}
\usepackage[T1]{fontenc}
\usepackage{lmodern,amsmath,amssymb,tabularx,booktabs,array}
\usepackage[table]{xcolor}
\usepackage[a4paper,left=22mm,right=22mm,top=22mm,bottom=21mm,headheight=26pt,headsep=12pt,footskip=25pt]{geometry}
\usepackage{fancyhdr,hyperref}
\input{glyphtounicode}
\pdfgentounicode=1
\definecolor{ink}{HTML}{132B3C}
\definecolor{teal}{HTML}{076E73}
\definecolor{muted}{HTML}{435D69}
\definecolor{rulegrey}{HTML}{B7CCCA}
\hypersetup{colorlinks=true,linkcolor=teal,urlcolor=teal,
  pdftitle={MH6003 Cryptography - Group project topics},
  pdfsubject={AI for cryptography, post-quantum cryptography and two further projects},pdfauthor={MH6003 Cryptography}}
\renewcommand{\familydefault}{\sfdefault}
\color{ink}
\setlength{\parindent}{0pt}
\setlength{\parskip}{4pt}
\setlength{\emergencystretch}{9pt}
\pagestyle{fancy}\fancyhf{}
\fancyhead[L]{\small\bfseries MH6003\enspace\textcolor{teal}{/}\enspace Cryptography}
\fancyhead[R]{\small Group project topics}
\fancyfoot[L]{\footnotesize\color{muted}Proposed topics for instructor review}
\fancyfoot[R]{\footnotesize\color{muted}\thepage}
\renewcommand{\headrule}{\hbox to\headwidth{\color{rulegrey}\leaders\hrule height .6pt\hfill}}
\newcommand{\pagetitle}[1]{{\fontsize{27}{30}\selectfont\bfseries #1\par}\vspace{5pt}}
\newcommand{\project}[3]{%
  \pdfbookmark[1]{#1. #2}{project#1}%
  \hypertarget{topic#1}{}{\small\bfseries\color{teal}TOPIC #1\enspace /\enspace #3\par}
  {\raggedright\fontsize{17}{20}\selectfont\bfseries #2\par}\vspace{1pt}}
\newcommand{\entry}[2]{{\raggedright\textcolor{teal}{\bfseries #1}\enspace #2\par}}
\newcommand{\reading}[1]{{\small\color{muted}Starting point\enspace #1\par}}
\newcommand{\topicbreak}{\vspace{6pt}{\color{rulegrey}\hrule}\vspace{7pt}}
\newcommand{\spot}[3]{\hyperlink{topic#1}{\textbf{#1}\enspace #2} & #3\\}
\newcommand{\track}[1]{\multicolumn{2}{@{}l}{\textcolor{teal}{\bfseries #1}}\\}
\begin{document}
\pdfbookmark[0]{Project menu}{menu}
\pagetitle{Group project topics}
{\large Three AI projects, three post-quantum projects and two further options.\par}
\medskip
\begin{tabularx}{\linewidth}{@{}*{3}{>{\raggedright\arraybackslash}X}@{}}
\toprule
\textbf{By week 7} & \textbf{Weeks 8--11} & \textbf{End of week 11}\\
Form a team of 3--4 and choose a topic. & Carry out the four-week project. & Submit the project report.\\
\bottomrule
\end{tabularx}

The group project contributes \textbf{30\% of the course grade}. A prototype or computational experiment is preferred. In the report, explain the research question, method, results and limitations. Identify each member's contribution.

\vspace{5pt}
{\large\bfseries Which question would your group like to answer?\par}
{\small Select a title to open its brief. Groups may also propose another cryptography question of comparable scope.\par}
\begingroup
\renewcommand{\arraystretch}{1.36}
\begin{tabularx}{\linewidth}{@{}>{\raggedright\arraybackslash}p{.48\linewidth} >{\raggedright\arraybackslash}X@{}}
\toprule
\textbf{Project} & \textbf{What the report could show}\\
\midrule
\track{AI for cryptography}
\spot{1}{Recognising classical ciphers}{Whether a classifier improves on a statistical rule}
\spot{2}{Learning to select small S-boxes}{Whether a model saves work in a component-design search}
\spot{3}{Checking an AI's reasoning about cipher rounds}{Which claims and counterexamples survive an exact check}
\midrule
\track{Post-quantum cryptography}
\spot{4}{The cost of post-quantum key establishment}{Bytes exchanged and time taken by each method}
\spot{5}{Post-quantum signatures for software updates}{How lattice and hash-based signatures affect a device fleet}
\spot{6}{Polynomial multiplication inside ML-KEM}{Why the transform multiplies coefficients in pairs}
\midrule
\track{Hash functions and privacy}
\spot{7}{Merkle proofs for a file catalogue}{How verification costs grow with catalogue size}
\spot{8}{Counting private responses with two servers}{What each server learns and what the tally reveals}
\bottomrule
\end{tabularx}
\endgroup

\newpage
\pagetitle{AI for cryptography}
\project{1}{Can a model recognise a classical cipher?}{Supervised classification}
\entry{Rationale}{Block transposition preserves letter counts, while a repeating Vigenère key can change them. Can a classifier use additional patterns to distinguish these ciphers on short passages?}
\entry{Suggested implementation}{Encrypt passages of three lengths from public-domain books both ways, using fresh keys and identical formatting. Choose lengths that avoid padding. Compare logistic regression on letter and adjacent-pair frequencies with a threshold on the index of coincidence. Fit both methods on training data. Hold out entire books and new keys for testing.}
\entry{Expected outcome}{Give a confusion matrix and accuracy by length for each method. Remove the adjacent-pair features and retest to see what they contributed. Explain one difficult passage.}
\reading{Lecture 1 on coincidence. \href{https://scikit-learn.org/stable/modules/cross_validation.html}{Scikit-learn guide to grouped evaluation}}
\topicbreak
\project{2}{Can learning reduce the cost of choosing an S-box?}{Machine learning for component design}
\entry{Rationale}{Consider four-bit substitution boxes (S-boxes) that permute the 16 possible values. Designers want to avoid input differences that too often yield the same output difference. A model could help decide which candidate tables merit exact scoring.}
\entry{Suggested implementation}{For each nonzero input difference $a$, count the occurrences of each $S(x)\oplus S(x\oplus a)$ over all 16 inputs $x$. The largest count is the score to minimise. Use the 16 table entries as features and train a small tree ensemble to predict the score. Have it rank fresh candidates, then score its choices exactly. Compare with random selection from the same pool.}
\entry{Expected outcome}{Across several seeds, plot the best exact score found against the number of labelled tables used, including the training set. Report runtime with model fitting included. Explain what this one design criterion measures.}
\reading{\href{https://doc.sagemath.org/html/en/reference/cryptography/sage/crypto/sbox.html}{SageMath S-box reference and difference tables}}
\topicbreak
\project{3}{Can a language model justify a reversibility claim?}{AI-assisted mathematical reasoning}
\entry{Rationale}{An eight-bit Feistel round can be invertible even when its inner function is many-to-one. Small maps let us check every input and see whether a language model can justify its claim.}
\entry{Suggested implementation}{Generate a balanced set of invertible and many-to-one eight-bit maps from simple Feistel and substitution rules. Enumerate all 256 inputs to label each map. On held-out examples, compare a direct yes/no prompt with one that asks for an inverse description or two different inputs with the same output. Use the same model version for both.}
\entry{Expected outcome}{Report accuracy and false claims of invertibility. Check proposed collisions directly and proposed inverses on all inputs. Walk through one response and show whether its justification holds.}
\reading{Lecture 2, slides 40--57, on permutations, substitution and Feistel rounds.}

\newpage
\pagetitle{Post-quantum cryptography}
\project{4}{What does an ML-KEM exchange cost?}{Key establishment}
\entry{Rationale}{Ephemeral X25519 uses elliptic curves. ML-KEM is a post-quantum key-encapsulation mechanism. Both let two peers establish a shared secret, but the messages they exchange differ. What does that difference cost on a connection with limited bandwidth?}
\entry{Suggested implementation}{Run X25519 and ML-KEM-768 between two local processes using existing libraries and fresh exchange keys. Count both X25519 public keys. For ML-KEM, count the recipient's public key and the ciphertext sent in response. Check that both peers obtain the same secret within each method. Time key generation and the remaining operations separately.}
\entry{Expected outcome}{Draw both message flows. Report bytes and median operation times with variability. Estimate transmission time at a stated link speed, distinguishing it from measured computation. Identify the authentication and framing costs outside the experiment.}
\reading{\href{https://doi.org/10.6028/NIST.FIPS.203}{FIPS 203, ML-KEM}\enspace$\cdot$\enspace\href{https://www.rfc-editor.org/rfc/rfc7748.html}{RFC 7748, X25519}}
\topicbreak
\project{5}{Which signature suits a software update?}{Lattice and hash-based signatures}
\entry{Rationale}{A publisher signs an update once, but every device downloads and verifies its signature. A short signature and a fast verifier can therefore matter differently to the publisher and to the devices.}
\entry{Suggested implementation}{Compare lattice-based ML-DSA-65 with hash-based SLH-DSA-SHA2-192s using library implementations. Both target NIST security category 3. Sign the same release manifest with each, using separate keys. Verify the original manifest and check that an edited manifest is rejected. Measure signature and public-key sizes, signing time and verification time on the same machine.}
\entry{Expected outcome}{For one release and 1,000 devices with preinstalled public keys, calculate total signature bytes downloaded and separate one-time signing from per-device verification. Recommend a scheme for a stated bandwidth or compute budget.}
\reading{\href{https://doi.org/10.6028/NIST.FIPS.204}{FIPS 204, ML-DSA}\enspace$\cdot$\enspace\href{https://doi.org/10.6028/NIST.FIPS.205}{FIPS 205, SLH-DSA}}
\topicbreak
\project{6}{Why does ML-KEM multiply coefficients in pairs?}{Lattice-scheme arithmetic}
\entry{Rationale}{ML-KEM computes polynomial products modulo $X^{256}+1$ with coefficients modulo 3329. Its number-theoretic transform (NTT) breaks this product into 128 small products. Each acts on a pair of coefficients.}
\entry{Suggested implementation}{Write a direct polynomial multiplier as a reference. Then implement the FIPS 203 forward transform, pair multiplication and inverse transform in the same language. Check products on monomials and seeded random inputs. One revealing boundary case is $X^{255}X=-1$ in this ring.}
\entry{Expected outcome}{Trace one coefficient pair through multiplication. Use the roots of unity available modulo 3329 to explain why the transform stops at pairs. Compare arithmetic counts and time for the full product, including transformations in both directions.}
\reading{\href{https://doi.org/10.6028/NIST.FIPS.203}{FIPS 203, NTT and multiplication algorithms}}

\newpage
\pagetitle{Hash functions and privacy}
\project{7}{How much data proves membership in a catalogue?}{Merkle trees}
\entry{Rationale}{A device knows the trusted root hash of an ordered file catalogue. To verify a file digest at a given position, it could download the entire list. A Merkle tree lets it check that one entry using the hashes along a single path. How does the saving change with catalogue size?}
\entry{Suggested implementation}{Build a binary tree over ordered lists of file digests using SHA-256 and separate prefixes for leaves and internal nodes. Start with powers of two. Generate and verify proofs, keeping track of sibling order. Test a changed digest and an incorrect path. Compare with sending the full ordered list and checking the same root.}
\entry{Expected outcome}{Produce a verifier and plot proof bytes and verification hashes against catalogue size. Derive the sibling-hash cost for $2^h$ entries and check it against the measurements. Explain how the device obtains a trusted root in your scenario.}
\reading{\href{https://www.rfc-editor.org/rfc/rfc9162.html\#section-2.1}{RFC 9162, Merkle trees and inclusion proofs}}
\topicbreak
\project{8}{Can two servers tally private responses?}{Secret sharing}
\entry{Rationale}{A class wants an exact count of yes votes. Each counting server can add its inputs while seeing only a random share of each vote. What information becomes visible once the tally is published?}
\entry{Suggested implementation}{Simulate binary votes and two servers that follow the protocol and do not collude. Choose a modulus $q$ larger than the number of voters. For each vote $v$, send a fresh uniform share $r$ modulo $q$ to one server and $(v-r)\bmod q$ to the other. Add the server totals modulo $q$. Compare the results and traffic with a direct tally at several class sizes.}
\entry{Expected outcome}{Check exact totals, including all-zero and all-one classes. Report bytes per voter under a stated encoding. Explain why either server's individual share is uniform for both possible votes, and give a case where the released total reveals someone's response.}
\reading{\href{https://csrc.nist.gov/projects/pec}{NIST Privacy-Enhancing Cryptography project}}
\end{document}
